% SOURCE_AUTHORITY: sga5_fr_workpass.tex, lines 5139--5178 (LF-normalized unit).
% SOURCE_SHA256: 791F4EFFC5E02832D5D77ED03518C8156D6F07E4C8238B03545DB93D883FBB28
\subsubsection*{5.8}

Sean $P,Q$ unos $A$-bimódulos. Entonces $P\otimes_AQ$ y $Q\otimes_AP$ son $A$-bimódulos, y se tienen isomorfismos canónicos
\begin{equation}\tag{5.8.1}
(P\otimes_AQ)\otimes_{A^e}A\xrightarrow{\sim}P\otimes_{A^e}Q\xrightarrow{\sim}(Q\otimes_AP)\otimes_{A^e}A\xrightarrow{\sim}A\otimes_{A^e}(P\otimes_KQ)\otimes_{A^e}A,
\end{equation}
identificándose canónicamente cada uno de los objetos anteriores con el $K$-módulo cociente de $P\otimes_KQ$ por el submódulo generado por las secciones de la forma $axb\otimes y-x\otimes bya$, para $x\in P$, $y\in Q$, $a,b\in A$. Si $A$ es plano sobre $K$, estos isomorfismos se derivan en isomorfismos
\begin{equation}\tag{5.8.2}
(P\lten_AQ)\lten_{A^e}A\simeq P\lten_{A^e}Q\simeq(Q\lten_AP)\lten_{A^e}A\simeq A\lten_{A^e}(P\lten_KQ)\lten_{A^e}A
\end{equation}
para $P,Q\in\ob D^-(A^e)$. Por otra parte, la flecha de composición en el nivel de los módulos
\[
\uHom_A(E,P\otimes_AF)\otimes_K\uHom_A(F,Q\otimes_AE)
\longrightarrow\uHom_A(E,P\otimes_AQ\otimes_AE),
\]
\[
f\otimes g\longmapsto (P\otimes_Ag)f,
\]
se deriva en
\begin{equation}\tag{5.8.3}
\uRHom_A(E,P\lten_AF)\lten_K\uRHom_A(F,Q\lten_AE)
\longrightarrow\uRHom_A(E,P\lten_AQ\lten_AE)
\end{equation}
para $P,Q\in\ob D^b(A^e)$, $E,F\in\ob D^b(A)$ tales que los productos tensoriales y los $\uRHom$ que figuran arriba estén en $D^b$. Cuando $E$ y $F$ son perfectos, esta flecha se identifica, mediante 5.4.2, con la flecha deducida de la flecha de evaluación $F\lten_K\uRHom_A(F,Q)\to Q$; de ello se deduce fácilmente que se tiene un cuadrado conmutativo
\begin{equation}\tag{5.8.4}
\begin{tikzcd}[column sep=large,row sep=large]
\uRHom_A(E,P\lten_AF)\lten_K\uRHom_A(F,Q\lten_AE) \arrow[r,"(1)"] \arrow[d,"(2)"'] & \uRHom_A(E,P\lten_AQ\lten_AE) \arrow[d,"\Tr_A"]\\
\uRHom_A(F,Q\lten_AP\lten_AF) \arrow[r,"\Tr_A"'] & P\lten_{A^e}Q,
\end{tikzcd}
\end{equation}
donde (1) y (2) vienen dadas por 5.8.3 (el punto es que las dos flechas compuestas de 5.8.4 están definidas por el producto tensorial de las flechas de evaluación
\[
F\lten_K\uRHom_A(F,Q)\longrightarrow Q,\qquad
E\lten_K\uRHom_A(E,P)\longrightarrow P.
\]
De ello resulta que, para $u\in\uHom_A(E,P\lten_AF)$, $v\in\uHom_A(F,Q\lten_AE)$, se tiene
\begin{equation}\tag{5.8.5}
\Tr_A(vu)=\Tr_A(uv).
\end{equation}
A veces denotaremos por $\langle u,v\rangle_A$ el valor común de 5.8.5.
